
\documentclass{article}
\usepackage{colortbl}
\usepackage{makecell}
\usepackage{multirow}
\usepackage{supertabular}

\begin{document}

\newcounter{utterance}

\centering \large Interaction Transcript for game `referencegame', experiment `random\_grids', episode 16 with reflexion{-}llama{-}3.1{-}8b{-}instruct.
\vspace{24pt}

{ \footnotesize  \setcounter{utterance}{1}
\setlength{\tabcolsep}{0pt}
\begin{supertabular}{c@{$\;$}|p{.15\linewidth}@{}p{.15\linewidth}p{.15\linewidth}p{.15\linewidth}p{.15\linewidth}p{.15\linewidth}}
    \# & $\;$A & \multicolumn{4}{c}{Game Master} & $\;\:$B\\
    \hline

    \theutterance \stepcounter{utterance}  
    & & \multicolumn{4}{p{0.6\linewidth}}{
        \cellcolor[rgb]{0.9,0.9,0.9}{
            \makecell[{{p{\linewidth}}}]{
                \texttt{\tiny{[P1$\langle$GM]}}
                \texttt{You have attempted this same task before and did not succeed. These are the plans you wrote after those attempts.} \\
\texttt{Trial \#1: New plan: I failed to distinguish the target grid from the distractors because my description was too general. The target grid has a row with three X's in the second, third, and fourth columns, but so does Distractor grid 2. Next time, I will look for a more unique feature that is present in the target grid but not in the distractors. I will examine the rows and columns more carefully to find a distinctive pattern or characteristic that can be used to describe the target grid.} \\
\\ 
\texttt{Trial \#2: New plan: In my previous attempt, I focused on a row with three X's, but it turned out that this feature was not unique to the target grid. I should have looked more closely at the columns instead. Upon re{-}examining the grids, I notice that the target grid has a column with two X's and a □ in the third row, which is not present in the distractor grids. Next time, I will focus on describing this unique column feature to distinguish the target grid from the distractors.} \\
\\ 
\texttt{Trial \#3: New plan: In my previous attempts, I focused on rows and columns with specific patterns of X's and □'s, but these features were not unique to the target grid. I should have looked more closely at the overall structure of the grids and the relationships between the X's and □'s. Upon re{-}examining the grids, I notice that the target grid has a unique combination of X's and □'s that is not present in the distractor grids. Specifically, the target grid has a "V" shape formed by X's in the top{-}left and bottom{-}right corners, which is not present in the distractor grids. Next time, I will focus on describing this unique "V" shape to distinguish the target grid from the distractors.} \\
\\ 
\\ 
\texttt{You are given three grids, where each of them is 5 by 5 in size.} \\
\texttt{Grids have empty cells marked with "▢" and filled cells marked with "X".} \\
\texttt{Your task is to generate a referring expression that best describes the target grid while distinguishing it from the two other distractor grids.} \\
\texttt{The first grid is the target grid, and the following two grids are the distractors.} \\
\\ 
\texttt{Target grid:} \\
\\ 
\texttt{X X □ X □} \\
\texttt{□ X X □ X} \\
\texttt{X □ □ X □} \\
\texttt{□ X □ □ X} \\
\texttt{X □ X □ X} \\
\\ 
\texttt{Distractor grid 1:} \\
\\ 
\texttt{X X □ □ X} \\
\texttt{□ X X □ X} \\
\texttt{X □ X X □} \\
\texttt{□ X □ X X} \\
\texttt{X □ □ X X} \\
\\ 
\texttt{Distractor grid 2:} \\
\\ 
\texttt{X □ X X □} \\
\texttt{□ X □ X X} \\
\texttt{X X □ □ X} \\
\texttt{□ X X □ X} \\
\texttt{X □ X □ □} \\
\\ 
\texttt{Instruction: Describe the target grid.} \\
\texttt{Generate the referring expression starting with the tag "Expression: " for the given target grid. Omit any other text.} \\
            }
        }
    }
    & \\ \\

    \theutterance \stepcounter{utterance}  
    & \multicolumn{4}{p{0.6\linewidth}}{
        \cellcolor[rgb]{0.9,0.9,0.9}{
            \makecell[{{p{\linewidth}}}]{
                \texttt{\tiny{[P1$\rangle$GM]}}
                \texttt{Expression: The grid that has a row with three X's in the second, third, and fourth columns and a column with two X's and a □ in the third row.} \\
            }
        }
    }
    & & \\ \\

    \theutterance \stepcounter{utterance}  
    & & & \multicolumn{2}{p{0.3\linewidth}}{
        \cellcolor[rgb]{0.9,0.9,0.9}{
            \makecell[{{p{\linewidth}}}]{
                \texttt{\tiny{[GM$|$GM]}}
                \texttt{Expression: The grid that has a row with three X's in the second, third, and fourth columns and a column with two X's and a □ in the third row.} \\
            }
        }
    }
    & & \\ \\

    \theutterance \stepcounter{utterance}  
    & & \multicolumn{4}{p{0.6\linewidth}}{
        \cellcolor[rgb]{0.9,0.9,0.9}{
            \makecell[{{p{\linewidth}}}]{
                \texttt{\tiny{[GM$\rangle$P2]}}
                \texttt{You have attempted this same task before and did not succeed. These are the plans you wrote after those attempts.} \\
\texttt{Trial \#1: New plan: I failed to identify the correct grid because I didn't carefully examine the expression. I should have looked for the row with three X's in the second, third, and fourth columns in each grid. In my next attempt, I will carefully examine each grid and select the one that matches the given expression. I will look for the row with three X's in the second, third, and fourth columns in each grid and choose the grid that satisfies this condition.} \\
\\ 
\texttt{Trial \#2: New plan: I failed to identify the correct grid because I didn't carefully examine the expression. However, in my previous attempt, I only looked for the row with three X's in the first, second, and third columns, but the expression actually mentions the first, second, and third columns. What I missed was that the expression does not specify the order of the columns, it only mentions the columns. In my next attempt, I will carefully examine each grid and select the one that matches the given expression. I will look for the row with three X's in any three consecutive columns in each grid and choose the grid that satisfies this condition.} \\
\\ 
\texttt{Trial \#3: New plan: I failed to identify the correct grid because I didn't carefully examine the expression. However, in my previous attempts, I only looked for the row with three X's in the second, third, and fourth columns, but the expression actually mentions a column with two X's and a □ in the third row. What I missed was that the expression does not specify the exact position of the two X's and the □ in the column, it only mentions the row. In my next attempt, I will carefully examine each grid and select the one that matches the given expression. I will look for the row with three X's in any three consecutive columns in each grid and also examine each column in the third row to find the one that has two X's and a □, and choose the grid that satisfies both conditions.} \\
\\ 
\\ 
\texttt{You are given three grids, where each of them is 5 by 5 in size.} \\
\texttt{Grids have empty cells marked with "▢" and filled cells marked with "X".} \\
\texttt{You are also given a referring expression that describes one of the given grids.} \\
\texttt{Your task is to select the grid that matches the given referring expression.} \\
\texttt{Generate only the number (in text) of the grid that matches the given expression by selecting first, second, or third.} \\
\\ 
\texttt{First:} \\
\\ 
\texttt{X X □ □ X} \\
\texttt{□ X X □ X} \\
\texttt{X □ X X □} \\
\texttt{□ X □ X X} \\
\texttt{X □ □ X X} \\
\\ 
\texttt{Second:} \\
\\ 
\texttt{X X □ X □} \\
\texttt{□ X X □ X} \\
\texttt{X □ □ X □} \\
\texttt{□ X □ □ X} \\
\texttt{X □ X □ X} \\
\\ 
\texttt{Third:} \\
\\ 
\texttt{X □ X X □} \\
\texttt{□ X □ X X} \\
\texttt{X X □ □ X} \\
\texttt{□ X X □ X} \\
\texttt{X □ X □ □} \\
\\ 
\texttt{Expression: The grid that has a row with three X's in the second, third, and fourth columns and a column with two X's and a □ in the third row.} \\
\texttt{Question: Which grid does the expression refer to?} \\
\texttt{Start with the tag "Answer: ", followed by your selection. Omit any other text.} \\
            }
        }
    }
    & \\ \\

    \theutterance \stepcounter{utterance}  
    & & & \multicolumn{4}{p{0.6\linewidth}}{
        \cellcolor[rgb]{0.9,0.9,0.9}{
            \makecell[{{p{\linewidth}}}]{
                \texttt{\tiny{[GM$\langle$P2]}}
                \texttt{Answer: 1} \\
            }
        }
    }
     \\ \\

    \theutterance \stepcounter{utterance}  
    & & & \multicolumn{2}{p{0.3\linewidth}}{
        \cellcolor[rgb]{0.9,0.9,0.9}{
            \makecell[{{p{\linewidth}}}]{
                \texttt{\tiny{[GM$|$GM]}}
                \texttt{1} \\
            }
        }
    }
    & & \\ \\

    \theutterance \stepcounter{utterance}  
    & \multicolumn{4}{p{0.6\linewidth}}{
        \cellcolor[rgb]{0.95,0.95,0.95}{
            \makecell[{{p{\linewidth}}}]{
                \texttt{\tiny{[P2$\rightarrow$P1]}}
                \texttt{Post{-}episode Reflexion (not sent to GM):} \\
\texttt{New plan: In my previous attempts, I focused on specific patterns of X's and □'s in rows and columns, but these features were not unique to the target grid. I should have looked more closely at the overall structure of the grids and the relationships between the X's and □'s. Upon re{-}examining the grids, I notice that the target grid has a unique combination of X's and □'s that is not present in the distractor grids. Specifically, the target grid has a "V" shape formed by X's in the top{-}left and bottom{-}right corners, which is not present in the distractor grids. However, I also notice that the target grid has a unique "L" shape formed by X's in the top{-}right and bottom{-}left corners, which is not present in the distractor grids. Next time, I will focus on describing this unique "L" shape to distinguish the target grid from the distractors.} \\
            }
        }
    }
    & & \\ \\

    \theutterance \stepcounter{utterance}  
    & & & \multicolumn{4}{p{0.6\linewidth}}{
        \cellcolor[rgb]{0.95,0.95,0.95}{
            \makecell[{{p{\linewidth}}}]{
                \texttt{\tiny{[P1$\rightarrow$P2]}}
                \texttt{Post{-}episode Reflexion (not sent to GM):} \\
\texttt{New plan: I failed to identify the correct grid because I didn't carefully examine the expression. However, in my previous attempts, I only looked for the row with three X's in the second, third, and fourth columns, but the expression actually mentions a column with two X's and a □ in the third row. What I missed was that the expression does not specify the exact position of the two X's and the □ in the column, it only mentions the row. In my next attempt, I will carefully examine each grid and select the one that matches the given expression. I will look for the row with three X's in any three consecutive columns in each grid and also examine each column in the third row to find the one that has two X's and a □, and choose the grid that satisfies both conditions.} \\
            }
        }
    }
     \\ \\

\end{supertabular}
}

\end{document}
